\documentclass[11pt,a4paper]{article}
\usepackage[utf8]{inputenc}
\usepackage[spanish]{babel}
\usepackage{amsmath, amssymb}
\usepackage{geometry}
\geometry{top=2.5cm, bottom=2.5cm, left=2.5cm, right=2.5cm}

\begin{document}

\section*{Piense, dialogue y comparta}

\begin{itemize}
    \item[\textbf{1.}] Usted trabaja en una estación de radar para la Guardia costera. Mientras todos los demás salieron a almorzar, oye una llamada de auxilio de un barco que se hunde. El barco está localizado a una distancia de $51.2 \text{ km}$ de la estación y una orientación de $36^\circ$ al noroeste. En su pantalla de radar, ve las siguientes localizaciones de otros cuatro barcos:
    \begin{center}
    \begin{tabular}{cllc}
    \hline
    Barco & \shortstack{Distancia desde \\ la estación (km)} & Orientación & \shortstack{Rapidez máxima \\ (km/h)} \\
    \hline
    1 & 36.1 & $42^\circ$ noroeste & 30.0 \\
    2 & 37.3 & $61^\circ$ noroeste & 38.0 \\
    3 & 10.2 & $36^\circ$ noroeste & 32.0 \\
    4 & 51.2 & $79^\circ$ noroeste & 45.0 \\
    \hline
    \end{tabular}
    \end{center}
    ¡Rápido!, ¿con cuál barco se contacta para ayudar al barco que se está hundiendo? ¿Qué barco llegará en el intervalo de tiempo más corto? Supongamos que cada barco aceleraría rápidamente a su velocidad máxima y luego mantiene esa velocidad constante en una línea recta durante todo el viaje hacia el barco que se está hundiendo.
    \vspace*{9cm}

    \item[\textbf{2.}] \textbf{ACTIVIDAD} En un mapa de papel de Estados Unidos, localice Memphis, Albuquerque y Chicago. Dibuje un vector desde Albuquerque a Memphis y otro vector de Memphis a Chicago. Usando la escala del mapa, determine las distancias en línea recta entre Albuquerque y Memphis, y entre Memphis y Chicago. Utilice un transportador para medir los ángulos de sus dos vectores respecto a la latitud y líneas de longitud. De esta información, determine la distancia en línea recta en millas entre Albuquerque y Chicago.
    \vspace*{8cm}
\end{itemize}

\newpage

\section*{Problemas}

\subsection*{SECCIÓN 3.1 Sistemas coordenados}

\begin{itemize}
    \item[\textbf{1.}] Dos puntos en el plano $xy$ tienen coordenadas cartesianas $(2.00, -4.00) \text{ m}$ y $(-3.00, 3.00) \text{ m}$. Determine: (a) la distancia entre estos puntos y (b) sus coordenadas polares.
    \vspace*{6cm}

    \item[\textbf{2.}] Dos puntos en un plano tienen coordenadas polares $(2.50 \text{ m}, 30.0^\circ)$ y $(3.80 \text{ m}, 120.0^\circ)$. Determine: (a) las coordenadas cartesianas de estos puntos y (b) la distancia entre ellos.
    \vspace*{6cm}

    \item[\textbf{3.}] Las coordenadas polares de un cierto punto son $(r = 4.30 \text{ cm}, \theta = 214^\circ)$. (a) Encuentre sus coordenadas cartesianas $x$ y $y$. Encuentre las coordenadas polares de los puntos con coordenadas cartesianas, (b) $(-x, y)$, (c) $(-2x, -2y)$ y (d) $(3x, -3y)$.
    \vspace*{7cm}

    \item[\textbf{4.}] Sean las coordenadas polares del punto $(x, y)$, $(r, \theta)$. Determine las coordenadas polares para los puntos (a) $(-x, y)$, (b) $(-2x, -2y)$ y (c) $(3x, -3y)$.
    \vspace*{6cm}
\end{itemize}

\newpage

\subsection*{SECCIÓN 3.2 Cantidades vectoriales y escalares}

\begin{itemize}
    \item[\textbf{5.}] ¿Por qué no es posible la siguiente situación? Una patinadora se desliza a lo largo de una trayectoria circular. Ella define cierto punto en el círculo como su origen. Posteriormente, pasa a través de un punto en el cual la distancia que ha viajado a lo largo de la trayectoria desde el origen es más pequeña que la magnitud de su vector de desplazamiento desde el origen.
    \vspace*{6cm}
\end{itemize}

\newpage

\subsection*{SECCIÓN 3.3 Aritmética vectorial básica}

\begin{itemize}
    \item[\textbf{6.}] El vector $\vec{A}$ tiene una magnitud de $29$ unidades y apunta en la dirección $y$ positiva. Cuando el vector $\vec{B}$ se suma al vector $\vec{A}$, el vector resultante $\vec{A} + \vec{B}$ apunta en dirección $y$ negativa con una magnitud de $14$ unidades. Encuentre la magnitud y la dirección de $\vec{B}$.
    \vspace*{6cm}

    \item[\textbf{7.}] Una fuerza $\vec{F}_1$ de $6.00$ unidades de magnitud actúa sobre un objeto en el origen en una dirección $30.0^\circ$ sobre el eje $x$ positivo (figura P3.7). Una segunda fuerza $\vec{F}_2$ de $5.00$ unidades de magnitud actúa sobre el objeto en la dirección del eje $y$ positivo. Encuentre gráficamente la magnitud y la dirección de la fuerza resultante $\vec{F}_1 + \vec{F}_2$.
    \vspace*{6cm}

    \item[\textbf{8.}] Tres desplazamientos son $\vec{A} = 200 \text{ m}$ al sur, $\vec{B} = 250 \text{ m}$ al oeste y $\vec{C} = 150 \text{ m}$ a $30.0^\circ$ al noreste. (a) Construya un diagrama separado para cada una de las siguientes posibles formas de sumar estos vectores: $\vec{R}_1 = \vec{A} + \vec{B} + \vec{C}$; $\vec{R}_2 = \vec{B} + \vec{C} + \vec{A}$; $\vec{R}_3 = \vec{C} + \vec{B} + \vec{A}$. (b) Explique qué puede concluir al comparar los diagramas.
    \vspace*{8cm}

    \item[\textbf{9.}] En la figura P3.9 se muestran dos vectores de desplazamiento $\vec{A}$ y $\vec{B}$, cuya magnitud es de $3.00 \text{ m}$ cada uno. La dirección del vector $\vec{A}$ es $\theta = 30.0^\circ$. Encuentre gráficamente (a) $\vec{A} + \vec{B}$, (b) $\vec{A} - \vec{B}$, (c) $\vec{B} - \vec{A}$ y (d) $\vec{A} - 2\vec{B}$. (Reporte todos los ángulos en sentido contrario a las manecillas del reloj desde el eje $x$ positivo.)
    \vspace*{8cm}

    \item[\textbf{10.}] Un carro de montaña rusa se mueve $200 \text{ pies}$ horizontalmente y luego se eleva $135 \text{ pies}$ en un ángulo de $30.0^\circ$ sobre la horizontal. A continuación, viaja $135 \text{ pies}$ a un ángulo de $40.0^\circ$ hacia abajo. ¿Cuál es su desplazamiento desde su punto de partida? Use técnicas gráficas.
    \vspace*{8cm}
\end{itemize}

\newpage

\subsection*{SECCIÓN 3.4 Componentes de un vector y vectores unitarios}

\begin{itemize}
    \item[\textbf{11.}] Una minivan viaja en línea recta al norte en el carril derecho de una autopista a $28.0 \text{ m/s}$. Un camper pasa a la minivan y luego cambia del carril izquierdo al derecho. Mientras lo hace, la trayectoria del camper sobre el camino es un desplazamiento recto a $8.50^\circ$ al noreste. Para evitar chocar con la minivan, la distancia norte-sur entre la defensa trasera del camper y la defensa delantera de la minivan no deben disminuir. (a) ¿El camper puede conducirse para satisfacer este requisito? (b) Explique su respuesta.
    \vspace*{6cm}

    \item[\textbf{12.}] Una persona camina $25.0^\circ$ al noreste durante $3.10 \text{ km}$. ¿Qué distancia tendría que caminar hacia el norte y hacia el este para llegar a la misma posición?
    \vspace*{6cm}

    \item[\textbf{13.}] Su perro está corriendo entre la hierba en el patio trasero. Realiza desplazamientos sucesivos $3.50 \text{ m}$ al sur, $8.20 \text{ m}$ al noreste y $15.0 \text{ m}$ al oeste. ¿Cuál es el desplazamiento resultante?
    \vspace*{6cm}

    \item[\textbf{14.}] Dados los vectores $\vec{A} = 2.00\hat{i} + 6.00\hat{j}$ y $\vec{B} = 3.00\hat{i} - 2.00\hat{j}$, (a) dibuje la suma vectorial $\vec{C} = \vec{A} + \vec{B}$ y la diferencia vectorial $\vec{D} = \vec{A} - \vec{B}$. (b) Calcule $\vec{C}$ y $\vec{D}$, en términos de vectores unitarios. (c) Calcule $\vec{C}$ y $\vec{D}$ en términos de coordenadas polares, con ángulos medidos respecto del eje positivo $x$.
    \vspace*{7cm}
\end{itemize}

\newpage

\begin{itemize}
    \item[\textbf{15.}] La vista desde el helicóptero en la figura P3.15 muestra a dos personas jalando una mula terca. La persona de la derecha tira con una fuerza $\vec{F}_1$ de magnitud $120 \text{ N}$ y la dirección de $\theta_1 = 60.0^\circ$. La persona de la izquierda jala con una fuerza $\vec{F}_2$ de $80.0 \text{ N}$ y dirección de $\theta_2 = 75.0^\circ$. Encuentre (a) la fuerza única que es equivalente a las dos fuerzas que se muestran y (b) la fuerza que una tercera persona tendría que ejercer sobre la mula para hacer que la fuerza resultante sea igual a cero. Las fuerzas se miden en unidades de newtons (representados por N).
    \vspace*{6cm}

    \item[\textbf{16.}] Una pendiente de esquiar cubierta de nieve forma un ángulo de $35.0^\circ$ con la horizontal. Cuando un esquiador cae a plomo por la colina, una porción de nieve salpicada se proyecta con un desplazamiento máximo de $1.50 \text{ m}$ a $16.0^\circ$ de la vertical hacia arriba de la colina, como se muestra en la figura P3.16. Encuentre las componentes de su posición máxima (a) paralela a la superficie y (b) perpendicular a la superficie.
    \vspace*{6cm}

    \item[\textbf{17.}] Considere los tres vectores desplazamiento $\vec{A} = (3\hat{i} - 3\hat{j}) \text{ m}$, $\vec{B} = (\hat{i} - 4\hat{j}) \text{ m}$ y $\vec{C} = (-2\hat{i} + 5\hat{j}) \text{ m}$. Use el método de componentes para determinar (a) la magnitud y dirección del vector $\vec{D} = \vec{A} + \vec{B} + \vec{C}$ y (b) la magnitud y dirección de $\vec{E} = -\vec{A} - \vec{B} + \vec{C}$.
    \vspace*{6cm}

    \item[\textbf{18.}] El vector $\vec{A}$ tiene componentes $x$ y $y$ de $-8.70 \text{ cm}$ y $15.0 \text{ cm}$, respectivamente; el vector $\vec{B}$ tiene componentes $x$ y $y$ de $13.2 \text{ cm}$ y $-6.60 \text{ cm}$, respectivamente. Si $\vec{A} - \vec{B} + 3\vec{C} = 0$, ¿cuáles son los componentes de $\vec{C}$?
    \vspace*{6cm}
\end{itemize}

\newpage

\begin{itemize}
    \item[\textbf{19.}] El vector $\vec{A}$ tiene componentes $x, y$ y $z$ de $8.00, 12.0$ y $-4.00$ unidades, respectivamente. (a) Escriba una expresión vectorial para $\vec{A}$ en notación de vector unitario. (b) Obtenga una expresión en vectores unitarios para un vector $\vec{B}$ de un cuarto de longitud de $\vec{A}$ que apunte en la misma dirección que $\vec{A}$. (c) Obtenga una expresión en vectores unitarios para un vector $\vec{C}$ de tres veces la longitud de $\vec{A}$ que apunte en la dirección opuesta a la dirección de $\vec{A}$.
    \vspace*{6cm}

    \item[\textbf{20.}] Dados los vectores desplazamiento $\vec{A} = (3\hat{i} - 4\hat{j} + 4\hat{k})$ y $\vec{B} = (2\hat{i} + 3\hat{j} - 7\hat{k})$, encuentre las magnitudes de los siguientes vectores y exprese cada uno en términos de sus componentes rectangulares. (a) $\vec{C} = \vec{A} + \vec{B}$; (b) $\vec{D} = 2\vec{A} - \vec{B}$.
    \vspace*{6cm}

    \item[\textbf{21.}] El vector $\vec{A}$ tiene un componente $x$ negativo de $3.00$ unidades de longitud y un componente $y$ positivo de $2.00$ unidades de longitud. (a) Determine una expresión para $\vec{A}$ en notación de vectores unitarios. (b) Determine la magnitud y dirección de $\vec{A}$. (c) ¿Qué vector $\vec{B}$, cuando se suma a $\vec{A}$, da un vector resultante sin componente $x$ y una componente $y$ negativa de $4.00$ unidades de longitud?
    \vspace*{6cm}

    \item[\textbf{22.}] En la figura P3.22 se muestran tres vectores desplazamiento de una pelota de croquet, donde $|\vec{A}| = 20.0$ unidades, $|\vec{B}| = 40.0$ unidades y $|\vec{C}| = 30.0$ unidades. Encuentren (a) la resultante en notación de vectores unitarios y (b) la magnitud y dirección del desplazamiento resultante.
    \vspace*{6cm}
\end{itemize}

\newpage

\begin{itemize}
    \item[\textbf{23.}] (a) Con $\vec{A} = (6.00\hat{i} - 8.00\hat{j})$ unidades, $\vec{B} = (-8.00\hat{i} + 3.00\hat{j})$ unidades y $\vec{C} = (26.0\hat{i} + 19.0\hat{j})$ unidades, determine $a$ y $b$ tales que $a\vec{A} + b\vec{B} + \vec{C} = 0$. (b) Un estudiante aprendió que una sola ecuación no se puede resolver para determinar valores para más de una incógnita en ella. ¿Cómo podría explicarle que tanto $a$ como $b$ se pueden determinar a partir de la ecuación que se usó en el inciso (a)?
    \vspace*{6cm}

    \item[\textbf{24.}] El vector $\vec{B}$ tiene componentes $x, y$ y $z$ de $4.00, 6.00$ y $3.00$ unidades, respectivamente. Calcule (a) la magnitud de $\vec{B}$ y (b) los ángulos que $\vec{B}$ forma con los ejes coordenados.
    \vspace*{6cm}

    \item[\textbf{25.}] Use el método de componentes para sumar los vectores $\vec{A}$ y $\vec{B}$ que se muestran en la figura P3.9. Ambos vectores tienen magnitudes de $3.00 \text{ m}$ y el vector $\vec{A}$ hace un ángulo de $\theta = 30.0^\circ$ con el eje $x$. Exprese la resultante $\vec{A} + \vec{B}$ en notación de vector unitario.
    \vspace*{6cm}

    \item[\textbf{26.}] Una chica que entrega periódicos cubre su ruta viajando $3.00$ cuadras al oeste, $4.00$ cuadras al norte, y luego $6.00$ cuadras al este. (a) ¿Cuál es su desplazamiento resultante? (b) ¿Cuál es la distancia total que viaja?
    \vspace*{6cm}
\end{itemize}

\newpage

\begin{itemize}
    \item[\textbf{27.}] Un hombre que empuja un trapeador por el suelo hace que experimente dos desplazamientos. El primero tiene una magnitud de $150 \text{ cm}$ y forma un ángulo de $120^\circ$ con el eje $x$ positivo. El desplazamiento resultante tiene una magnitud de $140 \text{ cm}$ y se dirige a un ángulo de $35.0^\circ$ con el eje $x$ positivo. Encuentre la magnitud y dirección del segundo desplazamiento.
    \vspace*{6cm}

    \item[\textbf{28.}] La figura P3.28 muestra proporciones típicas de las anatomías del varón (m) y de la mujer (f). Los desplazamientos $\vec{d}_{1m}$ y $\vec{d}_{1f}$ de la planta de los pies hasta el ombligo tienen magnitudes de $104 \text{ cm}$ y $84.0 \text{ cm}$, respectivamente. Los desplazamientos $\vec{d}_{2m}$ y $\vec{d}_{2f}$ desde el ombligo hasta las yemas de los dedos extendidas tienen magnitudes de $100 \text{ cm}$ y $86.0 \text{ cm}$, respectivamente. Encuentre la suma vectorial de estos desplazamientos $\vec{d}_3 = \vec{d}_1 + \vec{d}_2$ para las dos personas.
    \vspace*{6cm}

    \item[\textbf{29.}] \textbf{Problema de repaso.} Conforme pasa sobre la isla Gran Bahamas, el ojo de un huracán se mueve en una dirección $60.0^\circ$ al noroeste con una rapidez de $41.0 \text{ km/h}$. (a) ¿Cuál es la expresión en vector unitario para la velocidad del huracán? Mantiene esta velocidad por $3.00 \text{ h}$, después el curso del huracán cambia súbitamente al norte y su rapidez baja a $25.0 \text{ km/h}$. Esta nueva velocidad se mantiene durante $1.50 \text{ h}$. (b) ¿Cuál es la expresión en vector unitario de la nueva velocidad del huracán? (c) ¿Cuál es la expresión en vector unitario para el desplazamiento del huracán durante las primeras $3.00 \text{ h}$? (d) ¿Cuál es la expresión en vector unitario para el desplazamiento del huracán durante la última $1.50 \text{ h}$? (e) ¿A qué distancia de Gran Bahamas está el ojo $4.50 \text{ h}$ después de que pasa sobre la isla?
    \vspace*{8cm}
\end{itemize}

\newpage

\begin{itemize}
    \item[\textbf{30.}] En la operación de ensamblado que se ilustra en la figura P3.30, un robot primero mueve un objeto en línea recta hacia arriba y después también al este, alrededor de un arco que forma un cuarto de círculo de $4.80 \text{ cm}$ de radio que se encuentra en un plano vertical este-oeste. Luego el robot mueve el objeto hacia arriba y al norte, a través de un cuarto de círculo de $3.70 \text{ cm}$ de radio que se encuentra en un plano vertical norte-sur. Encuentre (a) la magnitud del desplazamiento total del objeto y (b) el ángulo que el desplazamiento total forma con la vertical.
    \vspace*{6cm}

    \item[\textbf{31.}] \textbf{Problema de repaso.} Usted está de pie sobre el suelo en el origen de un sistema coordenado. Un avión vuela sobre usted con velocidad constante paralela al eje $x$ y a una altura fija de $7.60 \times 10^3 \text{ m}$. En el tiempo $t = 0$ el avión está directamente arriba de usted de modo que el vector que va de usted a él es $\vec{P}_0 = 7.60 \times 10^3\hat{j} \text{ m}$. En $t = 30.0 \text{ s}$, el vector de posición que va de usted al avión es $\vec{P}_{30} = (8.04 \times 10^3\hat{i} + 7.60 \times 10^3\hat{j}) \text{ m}$ como se ve en la figura 3.31. Determine la magnitud y orientación del vector de posición del avión en $t = 45.0 \text{ s}$.
    \vspace*{6cm}

    \item[\textbf{32.}] ¿Por qué no es posible la siguiente situación? Un comprador que empuja un carrito por un mercado sigue las instrucciones hacia los productos enlatados y tiene un desplazamiento $8.00\hat{i} \text{ m}$ por un pasillo. Luego hace un giro de $90.0^\circ$ y se mueve $3.00 \text{ m}$ a lo largo del eje $y$. A continuación, hace otro giro de $90.0^\circ$ y se mueve $4.00 \text{ m}$ a lo largo del eje $x$. Cada comprador que siga estas instrucciones correctamente termina a $5.00 \text{ m}$ del punto de partida.
    \vspace*{6cm}

    \item[\textbf{33.}] En la figura P3.33, el segmento de recta representa una trayectoria desde el punto con vector de posición $(5\hat{i} + 3\hat{j}) \text{ m}$ al punto con posición $(16\hat{i} + 12\hat{j}) \text{ m}$. El punto A está en dicha trayectoria, a una fracción $f$ del camino hacia el destino. (a) Encuentre el vector de posición del punto A en términos de $f$. (b) Evalúe la expresión del inciso (a) en el caso $f = 0$. (c) Explique si el resultado en (b) es razonable. (d) Evalúe la expresión para $f = 1$. (e) Explique si el resultado en (d) es razonable.
    \vspace*{8cm}
\end{itemize}

\newpage

\subsection*{PROBLEMAS ADICIONALES}

\begin{itemize}
    \item[\textbf{34.}] Usted está pasando el verano como aprendiz para navegar en un barco grande que transporta carga por el lago Erie. Un día, usted y su barco deben viajar a través del lago a una distancia de $200 \text{ km}$ hacia el norte de su puerto de origen a su puerto de destino. Justo cuando sale de su puerto de origen, el sistema de navegación electrónica de la nave se descompone. El capitán continúa navegando, afirmando que puede depender de sus años de experiencia en el agua para guiarse. Los ingenieros trabajan en el sistema de navegación mientras el barco continúa avanzando, y los vientos y las olas lo empujan fuera de curso. Eventualmente, el sistema de navegación funciona lo suficiente para indicararle su ubicación. El sistema le indica que su posición actual es $50.0 \text{ km}$ al norte del puerto de origen y $25.0 \text{ km}$ al este del puerto. El capitán está un poco avergonzado de que su barco esté tan lejos de curso y le grita una orden a usted para que le diga de inmediato qué rumbo debe establecer desde su posición actual al puerto de destino. Dele un ángulo de rumbo apropiado.
    \vspace*{6cm}

    \item[\textbf{35.}] Una persona que sale a caminar sigue la trayectoria que se muestra en la figura P3.35. El viaje total consiste en cuatro trayectorias en línea recta. Al final de la caminata, ¿cuál es el desplazamiento resultante de la persona, medido desde el punto de partida?
    \vspace*{6cm}

    \item[\textbf{36.}] Un transbordador lleva turistas entre tres islas. Navega de la primera isla a la segunda, a $4.76 \text{ km}$ de distancia, en una dirección $37.0^\circ$ al noreste. Luego navega de la segunda isla a la tercera en una dirección de $69.0^\circ$ al noroeste. Por último, regresa a la primera isla y navega en una dirección $28.0^\circ$ al sureste. Calcule la distancia entre (a) la segunda y tercera islas y (b) la primera y tercera islas.
    \vspace*{6cm}

    \item[\textbf{37.}] Dos vectores $\vec{A}$ y $\vec{B}$ tienen magnitudes exactamente iguales. Para que la magnitud de $\vec{A} + \vec{B}$ sea $100$ veces más grande que la magnitud de $\vec{A} - \vec{B}$, ¿cuál debe ser el ángulo entre ellos?
    \vspace*{6cm}
\end{itemize}

\newpage

\begin{itemize}
    \item[\textbf{38.}] Dos vectores $\vec{A}$ y $\vec{B}$ tienen magnitudes exactamente iguales. Para que la magnitud de $\vec{A} + \vec{B}$ sea $n$ veces más grande que la magnitud de $\vec{A} - \vec{B}$, ¿cuál debe ser el ángulo entre ellos?
    \vspace*{6cm}

    \item[\textbf{39.}] \textbf{Problema de repaso.} El animal de peluche más grande del mundo es una víbora de $420 \text{ m}$ de largo, construida por niños noruegos. Suponga que la víbora se encuentra en un parque, como se muestra en la figura P3.39, y forma dos lados rectos de un ángulo de $105^\circ$, con un lado de $240 \text{ m}$ de largo. Olaf e Inge corren una competencia que inventan. Inge corre directamente desde la cola de la víbora a su cabeza, y Olaf parte del mismo lugar en el mismo momento, pero corre a lo largo de la víbora. (a) Si ambos niños corren uniformemente a $12.0 \text{ km/h}$, ¿cuánto tiempo antes que Olaf, Inge llega a la cabeza de la víbora? (b) Si Inge corre la carrera de nuevo con una rapidez constante de $12.0 \text{ km/h}$, ¿con qué rapidez constante debe correr Olaf para llegar al final de la serpiente al mismo tiempo que Inge?
    \vspace*{6cm}

    \item[\textbf{40.}] Ecoturistas utilizan su indicador de sistema de posicionamiento global para determinar su ubicación dentro de un jardín botánico como latitud $0.002\ 43$ grados al sur del ecuador, longitud $75.642\ 38$ grados oeste. Ellos desean visitar un árbol en la latitud $0.001\ 62$ grados norte, longitud $75.644\ 26$ grados oeste. (a) Determine la distancia de la línea recta y la dirección en la que se puede caminar para alcanzar el árbol de la siguiente manera. Primero modele la Tierra como una esfera de radio $6.37 \times 10^6 \text{ m}$ para determinar las componentes hacia el oeste y hacia el norte de desplazamiento necesario, en metros. Luego modele la Tierra como un plano superficial para completar el cálculo. (b) Explique por qué es posible utilizar estos dos modelos geométricos juntos para resolver el problema.
    \vspace*{6cm}

    \item[\textbf{41.}] Un vector está dado por $\vec{R} = 2\hat{i} + \hat{j} + 3\hat{k}$. Encuentre (a) las magnitudes de los componentes $x, y$ y $z$, (b) la magnitud de $\vec{R}$ y (c) los ángulos entre $\vec{R}$ y los ejes $x, y$ y $z$.
    \vspace*{6cm}
\end{itemize}

\newpage

\begin{itemize}
    \item[\textbf{42.}] Usted está trabajando como asistente de un controlador de tráfico aéreo en el aeropuerto local, del cual aterrizan y despegan los pequeños aeroplanos. Su trabajo es asegurarse de que los aviones no estén más cerca uno del otro que una distancia de separación segura mínima de $2.00 \text{ km}$. Usted observa dos pequeñas aeronaves en su pantalla de radar, sobre la superficie del océano. La primera está a una altitud $800 \text{ m}$ sobre la superficie, distancia horizontal de $19.2 \text{ kilómetros}$ y $25.0^\circ$ al suroeste. El segundo avión está a una altitud $1\,100 \text{ m}$, a una distancia horizontal de $17.6 \text{ km}$, y $20.0^\circ$ suroeste. A su supervisor le preocupa que los dos aviones estén demasiado cerca y le pide la distancia de separación de los dos aviones. (Coloque el eje $x$ al oeste, el eje $y$ al sur y el eje $z$ vertical.)
    \vspace*{7cm}

    \item[\textbf{43.}] \textbf{Problema de repaso.} La posición instantánea de un objeto se especifica por su vector de posición dirigido desde un origen fijo a la posición del objeto, representado como partícula. Suponga para cierto objeto que el vector de posición es una función de tiempo dado por $\vec{r} = 4\hat{i} + 3\hat{j} - 2t\hat{k}$, donde $\vec{r}$ está en metros y $t$ en segundos. (a) Evalúe $d\vec{r}/dt$. (b) ¿Qué cantidad física representa $d\vec{r}/dt$ para el objeto?
    \vspace*{6cm}

    \item[\textbf{44.}] Los vectores $\vec{A}$ y $\vec{B}$ tienen magnitudes iguales de $5.00$. La suma de $\vec{A}$ y $\vec{B}$ es el vector $6.00\hat{j}$. Determine el ángulo entre $\vec{A}$ y $\vec{B}$.
    \vspace*{6cm}

    \item[\textbf{45.}] Un paralelepípedo rectangular tiene dimensiones $a, b$ y $c$ como se muestra en la figura P3.45. (a) Obtenga una expresión vectorial para el vector diagonal de la cara $\vec{R}_1$. (b) ¿Cuál es la magnitud de este vector? (c) Observe que $\vec{R}_1$, $c\hat{k}$ y $\vec{R}_2$ forman un triángulo rectángulo. Obtenga una expresión vectorial para el vector diagonal cuerpo $\vec{R}_2$.
    \vspace*{6cm}
\end{itemize}

\newpage

\subsection*{PROBLEMA DE DESAFÍO}

\begin{itemize}
    \item[\textbf{46.}] Un pirata ha enterrado su tesoro en una isla con cinco árboles situados en los puntos $(30.0 \text{ m}, -20.0 \text{ m})$, $(60.0 \text{ m}, 80.0 \text{ m})$, $(-10.0 \text{ m}, -10.0 \text{ m})$, $(40.0 \text{ m}, -30.0 \text{ m})$ y $(-70.0 \text{ m}, 60.0 \text{ m})$, todos medidos respecto a algún origen, como se muestra en la figura P3.46. Su bitácora le instruye comenzar en el árbol $A$ y avanzar hacia el árbol $B$, pero cubrir solo la mitad de la distancia entre $A$ y $B$. Luego, moverse al árbol $C$, cubrir un tercio de la distancia entre su ubicación actual y $C$. Después avanzar hacia el árbol $D$, cubriendo un cuarto de la distancia entre donde está y $D$. Finalmente avanzar hacia el árbol $E$, cubriendo una quinta parte de la distancia entre usted y $E$, detenerse y cavar. (a) Suponga que tiene correctamente determinando el orden en que el pirata etiquetó los árboles como $A, B, C, D$ y $E$ como se muestra en la figura. ¿Cuáles son las coordenadas del punto donde dejó su tesoro enterrado? (b) \textbf{¿Qué pasaría si?} ¿Y si no se sabe muy bien la forma en que el pirata etiqueta los árboles? ¿Cómo modificaría la respuesta si ha cambiado el orden de los árboles, por ejemplo, a $B(30 \text{ m}, -20 \text{ m})$, $A(60 \text{ m}, 80 \text{ m})$, $E(-10 \text{ m}, -10 \text{ m})$, $C(40 \text{ m}, -30 \text{ m})$ y $D(-70 \text{ m}, 60 \text{ m})$? Establezca un razonamiento para mostrar que la respuesta no depende del orden en que se etiquetan los árboles.
    \vspace*{9cm}
\end{itemize}

\end{document}
